\documentclass[11pt,a4paper]{article}
\usepackage[margin=2.5cm]{geometry}
\usepackage{amsmath,amssymb}
\usepackage{booktabs}
\usepackage{parskip}
\usepackage{hyperref}
\usepackage[table]{xcolor}
\usepackage{graphicx}
\usepackage{colortbl}
\usepackage{siunitx}
\usepackage{verbatim}
\usepackage{subcaption}
\graphicspath{{figures_2026-06-02/}}

\hypersetup{colorlinks=true,linkcolor=blue,citecolor=blue,urlcolor=blue}

\newcommand{\DS}{\Delta\Sigma}
\newcommand{\dm}{\Delta m}
\newcommand{\Msun}{\,M_{\odot}}
\newcommand{\Mpc}{\,\mathrm{Mpc}}
\newcommand{\Mtwoh}{M_{200m}}

\title{\textbf{From the WaZP Cluster Catalog to a Mass--Richness Relation}\\[0.4em]
\large End-to-end weak-lensing pipeline and first results}
\author{destacking project}
\date{2026 June 2}

\begin{document}
\maketitle

\begin{abstract}
We document the construction and first scientific outputs of a complete
weak-lensing pipeline applied to the DES Y6 WaZP cluster catalog with DES Y3
metacalibration shape estimates. The pipeline (i) selects a clean WaZP
sub-sample of $15{,}967$ clusters, (ii) measures per-cluster magnitude gaps
in $i$-band, (iii) measures per-cluster $\DS(R)$ profiles using a fully
parallelized metacalibration cone-query, (iv) stacks the signal in three
dimensions of $(z, N_{\rm gals}, \dm)$, (v) fits each stack with a full
miscentered NFW model via nested sampling, and (vi) extracts a power-law
mass--richness relation with redshift evolution. A companion analysis splits
each $(z,N_{\rm gals})$ cell at the cell median magnitude gap to test whether
the $M$--$\lambda$ relation depends on $\dm$. We summarize methodology,
key numerical results, caveats, and outstanding issues.
\end{abstract}

\tableofcontents
\bigskip

\section{Overview of the pipeline}

The work is organized as five notebooks, each producing an artifact that is
consumed by the next stage. \autoref{tab:flow} lists them in execution order
and their outputs. All paths are relative to the repository root.

\begin{table}[h]
\centering
\small
\begin{tabular}{p{0.32\linewidth}p{0.30\linewidth}p{0.30\linewidth}}
\toprule
Notebook & Reads & Produces \\
\midrule
\verb|wazp_cluster_selection| & WaZP cluster + member FITS &
\verb|data/wazp_real_data.fits| \\
\verb|wazp_lensing_signals| & WaZP real data + DES Y3 metacal HDF5 + DNF photo-z HDF5 &
\verb|data/metacal_filtered.npz|,
\verb|data/wazp_lensing_signals.pkl| \\
\verb|wazp_stacked_signals| & metacal pre-filter + WaZP real data &
\verb|data/wazp_bkgs.pkl|,
\verb|data/wazp_stacked_signals.pkl| \\
\verb|wazp_mass_richness| & \verb|wazp_bkgs.pkl| + WaZP real data &
\verb|data/wazp_mass_richness_*.pkl| \\
\verb|wazp_mass_richness_by_delta_m| & idem &
\verb|data/wazp_mr_dm_*.pkl| \\
\bottomrule
\end{tabular}
\caption{Notebook execution order, inputs, and persisted outputs.}
\label{tab:flow}
\end{table}

The pipeline is reproducible from the WaZP and DES~Y3 raw HDF5/FITS files
with no manual steps; only the cosmology choice (\textsc{astropy}
\verb|Planck18|) is hard-coded.

\section{Sample selection (\texttt{wazp\_cluster\_selection})}
\label{sec:selection}

\paragraph{Inputs.}
\begin{itemize}
\item \verb|y6a2_dnf_wazp_v5.0.12.6801_clusters.fits|: 416{,}947 WaZP clusters,
\verb|ZPHOT| in $0.027$--$1.317$, \verb|NGALS| in $5.0$--$308.8$.
\item \verb|y6a2_dnf_wazp_v5.0.12.6801_members.fits|: 13{,}133{,}244 member rows
with columns including \verb|ID_CLUSTER|, \verb|FLAG_CG|, \verb|PMEM|,
\verb|MAG_I|, sky position.
\end{itemize}

\paragraph{Projected distance per member.}
For every member we compute the projected proper distance to its parent cluster's
centre,
\[
  R = \theta(\mathrm{RA},\mathrm{DEC}) \cdot D_A(z_{\rm cl}),
\]
with $\theta$ the haversine angular separation between member and cluster
position and $D_A$ from \verb|Planck18|. $D_A$ is computed once per cluster
and broadcast to its members. We use the WaZP detection centre
(\verb|RA|, \verb|DEC|) since \verb|RA_CG|/\verb|DEC_CG| are NaN-only in the
release. Sanity checks confirm $97.4\%$ of members fall inside their parent
\verb|RADIUS_MPC|; BCGs have a median $R_{\rm MPC}=0.085\Mpc$.

\paragraph{$i$-band magnitude gap.}
For each cluster we compute the expected magnitude gap between the BCG and
the brightest satellite within $r_{\rm cut}=0.5\Mpc$, using the
\verb|destacking.data.magnitude_gaps| function (which we fixed; the original
contained syntax errors and several variable typos that prevented it from
ever running). The function weights every candidate ``brightest satellite''
by the probability that it is a member \emph{and} every brighter satellite is
not, so $\dm$ is marginalized over membership uncertainty via \verb|PMEM|.
Result: $336{,}379$ clusters have finite $\dm_i$, with median $+0.517$~mag,
mean $+0.620$~mag, and $16$/$84$-percentile range
$[+0.140,\,+1.133]\,\mathrm{mag}$. The negative tail extends to
$-2.2\,\mathrm{mag}$ and is most likely a cluster-finder artefact
(BCG mis-assignment); the high tail extends to $+4.2\,\mathrm{mag}$.

\paragraph{Selection cuts.}
We apply three cuts \emph{after} computing $\dm$ so the analysis catalog
carries the gap as a column:
\begin{align*}
  N_{\rm gals} &> 25, \\
  z_{\rm phot} &< 0.7, \\
  \text{at least one member with \texttt{FLAG\_CG = True}}.
\end{align*}
The cuts survive 10.15\%, 57.08\%, and 80.68\% of the catalog individually
and \emph{15{,}967 clusters} when intersected
(see~\autoref{tab:cuts}). All 15{,}967 selected clusters have a
finite $\dm_i$. The final catalog \verb|data/wazp_real_data.fits| carries
\verb|(ID, RA, DEC, NGALS, ZPHOT, DELTA_M_I, DELTA_M_I_ERR)|.

\begin{figure}[h]
\centering
\includegraphics[width=0.85\linewidth]{rmpc_dist.png}
\caption{Projected proper-distance distribution of all WaZP members and
their distribution relative to the parent cluster's WaZP detection radius.
Members are heavily concentrated to small $R$, and $97.4\%$ lie inside the
catalog \texttt{RADIUS\_MPC}.}
\label{fig:rmpc}
\end{figure}

\begin{figure}[h]
\centering
\includegraphics[width=0.92\linewidth]{magnitude_gaps.png}
\caption{$i$-band magnitude-gap statistics for the WaZP catalog. Top-left:
distribution of $\dm_i$ over $336{,}379$ clusters with a BCG and at least
one inner satellite; median $+0.52$\,mag. Top-right: per-cluster
$\sigma(\dm_i)$ from the probabilistic estimator. Bottom row: $\dm_i$ vs
log richness and vs photo-$z$, showing a mild positive correlation with
richness and no strong $z$-trend.}
\label{fig:dm}
\end{figure}

\begin{table}[h]
\centering
\small
\begin{tabular}{lrr}
\toprule
Cut & Surviving & Fraction \\
\midrule
Total clusters                      & 416{,}947 & 100.00\,\% \\
\verb|NGALS > 25|                   &  42{,}322 &  10.15\,\% \\
\verb|ZPHOT < 0.7|                  & 237{,}975 &  57.08\,\% \\
Has BCG (\verb|FLAG_CG=True| in members) & 336{,}379 &  80.68\,\% \\
\midrule
All three cuts                      & 15{,}967  &   3.83\,\% \\
\bottomrule
\end{tabular}
\caption{Per-cut survival of the WaZP real-data selection.}
\label{tab:cuts}
\end{table}

\section{Per-cluster lensing signals (\texttt{wazp\_lensing\_signals})}
\label{sec:percluster}

\paragraph{Source pre-filter.} The DES Y3 metacalibration shape file
(\SI{291}{GB}, $399{,}263{,}026$ rows) is streamed in $10^7$-row chunks
together with the DES Y3 GOLD DNF photo-z file
(\verb|DESY3_GOLD_2_2.1_DNF.h5|, also $399{,}263{,}026$ rows; row-aligned
with metacal as verified by the \verb|coadd_object_id| at first/last
indices). Standard metacal cuts
($\texttt{flags}=0$, $\texttt{snr}\in[10,1000]$,
$\texttt{size\_ratio}>0.5$, $T>0$) plus $\texttt{zmean\_sof}>0$ leave
$126{,}719{,}902$ sources ($31.7\%$ retention). Only the 10 columns
required by \verb|xlensing.data.metacal_cluster_lensing|
(\verb|RA, DEC| in radians; $z$; $e_1, e_2$; weight; $R_{11}, R_{12}, R_{21},
R_{22}$) are kept, stored as \verb|float32| in \verb|data/metacal_filtered.npz|.
Wall time: $\sim 4.6\,\mathrm{min}$ on one read pass.

\paragraph{Spatial index.} We build a \verb|sklearn.neighbors.BallTree| over
$(\mathrm{dec}, \mathrm{ra})$ in radians with the haversine metric. Per-cluster
cone radius is $\theta_{\max} = R_{\rm Mpc}/D_A(z_{\rm cl})$ with
$R_{\rm Mpc}=5\Mpc$. Build time over $\sim 1.27 \times 10^8$ points:
\SI{\sim 80}{\second}.

\paragraph{Radial bin layout.} We use the
\verb|prototype_destacking| convention,
$\mathrm{radii} = \texttt{np.logspace}(-0.6,\,0.6,\,8)$ with edges offset by
half a log-bin width so the radii fall at bin centres. Bins span
$[0.21,\,4.85]\Mpc$.

\paragraph{Worker and pool.} A worker function performs one cone query,
constructs the 10-tuple in the order
\verb|metacal_cluster_lensing| expects, calls that function (which applies its
own $z_{\rm src}>1.1\,z_{\rm cl}+0.2$ background cut and computes
$\Sigma_{\rm crit}$), packages the output into the 6-row \verb|bkg| array
$\bigl[\Sigma_{\rm crit}, e_t, e_x, w, R, M\bigr]$, then calls
\verb|xlensing.data.single_cluster([bkg], BINS_LIMS, NBOOT=200)| to obtain
the bootstrap-resampled $\DS$ profile and covariance. The source arrays and
\verb|BallTree| are module-level globals so forked workers inherit them via
copy-on-write -- no per-task pickling cost. With $8$ workers on the dev
machine the $15{,}967$ profiles complete in $\sim 30$\,minutes.

\paragraph{Result.} $15{,}929$ usable profiles, $38$ empty (no sources in the
cone), $1$ error (a low-$z$ cluster with empty bins). Smoke test on cluster
\verb|ID 415758|: $2358$ background sources,
$\DS = [201, 272, 271, 169, 59, -12, 35, 7]\,M_\odot/{\rm pc}^2$ in
the eight radial bins, diagonal errors $\sim 30\text{--}320\,M_\odot/{\rm pc}^2$.

\begin{figure}[h]
\centering
\includegraphics[width=0.65\linewidth]{percluster_sample.png}
\caption{Per-cluster $\DS(R)$ scatter for $\sim 1000$ random clusters from
the WaZP real-data sample. Every point is one cluster in one radial bin.
The slope $\DS\propto R^{-1}$ visible in the inner bins (left side)
matches the expected NFW projection; the dispersion at fixed $R$ reflects
the mass distribution of the sample.}
\label{fig:percluster}
\end{figure}

\section{Stacked signals (\texttt{wazp\_stacked\_signals})}
\label{sec:stacked}

We construct a disjoint 3D nested partition of the sample into 27 stacks:
\begin{itemize}
\item \textbf{3 redshift slices} at fixed edges $[0.10,\,0.35,\,0.55,\,0.70]$;
\item within each slice, \textbf{3 richness slices} at the same fixed edges
$[25,\,40,\,70,\,\infty)$;
\item within each $(z,N_{\rm gals})$ cell, \textbf{3 $\dm$ tertiles}
computed \emph{per-cell} from finite $\dm$ values.
\end{itemize}
A cluster belongs to at most one of the 27 stacks. The selection survives
$15{,}913$ clusters ($54$ excluded for falling outside the $z/N_{\rm gals}$
edges or having non-finite $\dm$, although in our real-data file every cluster
already has finite $\dm$). Cell populations are large (smallest: $295$
clusters; largest: $4{,}940$), guaranteeing per-cell $\dm$ tertiles are well
defined. Each cell is stacked via \texttt{xlensing.data.stacked\_signal(bkgs, BINS\_LIMS, NBOOT=500)}, which returns
$(\Sigma, \mathrm{boost}, \mathrm{cov}_\Sigma, \xi, \mathrm{cov}_\xi)$.

The stacked profiles are visualized as a $3\times3$ panel grid (rows = $z$
slice, cols = $N_{\rm gals}$ slice) with the three $\dm$ tertile curves
overlaid in each panel. ${\DS}$ vs $R$ is plotted log--log.

\section{Mass--richness relation (\texttt{wazp\_mass\_richness})}
\label{sec:massrichness}

This notebook drops the $\dm$ split but increases the $N_{\rm gals}$
resolution to 6 log-spaced bins
$[25,\,37.9,\,57.4,\,87.1,\,132.2,\,200.0,\,\infty)$, giving 18 stacks
($3$ $z$ $\times$ $6$ $N_{\rm gals}$).

\paragraph{Cell population.} Six of the eighteen cells exceed $1000$
clusters; the smallest $(z\!=\![0.55,0.70),\,N_{\rm gals}\!>\!200)$
contains only \emph{two} clusters. ${\dm}$-aware analyses should treat the
two highest-$z$ highest-$\lambda$ stacks with caution.

\paragraph{Model.} We fit each stack with the full
\verb|xlensing.model.NFW_Delta_Sigma| signal --- centered NFW $+$ Rayleigh-
miscentered NFW $+$ BCG $+$ 2-halo --- with four free parameters and
$\texttt{BCG\_B\_MASS}=0$:
\begin{align*}
  \log_{10}(\Mtwoh / \Msun) &\in [13, 16], \\
  C_{200m} &\in [1, 15], \\
  F_{\rm miss} &\in [0, 1], \\
  \sigma_{\rm off} &\in [0.01, 1.0]\Mpc.
\end{align*}
The Gaussian likelihood uses the bootstrap covariance from
\verb|stacked_signal| with a small diagonal jitter
($10^{-6}\times\overline{\mathrm{diag(cov)}}$) for numerical conditioning.
The stack redshift is fixed to its mean. Posteriors are obtained with
\verb|nautilus.Sampler(n_live=400)|, run serially over the 18 stacks
(\SI{\sim 0.3}{\minute} each).

\paragraph{Best-fit masses.}
\autoref{tab:fits18} reports the weighted-mean posterior values. Twelve of
the eighteen stacks yield $\sim 4$ per cent $\log\Mtwoh$ uncertainties.
Concentration constraints are weak (most posteriors span $C\in[5,12]$
within the prior; uncertainty not shown). The $F_{\rm miss}$ posteriors
cluster around $0.3\text{--}0.5$ and $\sigma_{\rm off}\sim 0.2\text{--}0.6\Mpc$;
both are well-constrained only in the high-richness, low-$z$ cells. The
stack $\mathtt{z2\_ng5}$ ($N=2$) hits multiple prior edges
($C=14.99$, $F_{\rm miss}=0.01$, $\log\Mtwoh = 15.94$) and is treated as
an outlier.

\begin{table}[h]
\centering
\small
\begin{tabular}{lrrrrr}
\toprule
Stack & $\langle z \rangle$ & $\langle\lambda\rangle$ & $N$ &
$\log\Mtwoh$ & $\sigma_{\log M}$ \\
\midrule
z0\_ng0 & 0.262 &  30.3 & 1{,}448 & 14.232 & 0.028 \\
z0\_ng1 & 0.271 &  45.9 &   689 & 14.437 & 0.028 \\
z0\_ng2 & 0.272 &  69.3 &   329 & 14.647 & 0.026 \\
z0\_ng3 & 0.270 & 102.9 &   101 & 14.778 & 0.037 \\
z0\_ng4 & 0.284 & 155.4 &    38 & 15.049 & 0.041 \\
z0\_ng5 & 0.284 & 239.2 &     6 & 15.086 & 0.052 \\
z1\_ng0 & 0.451 &  30.2 & 4{,}556 & 14.155 & 0.022 \\
z1\_ng1 & 0.449 &  45.7 & 2{,}009 & 14.403 & 0.032 \\
z1\_ng2 & 0.446 &  68.6 &   867 & 14.522 & 0.028 \\
z1\_ng3 & 0.439 & 103.6 &   266 & 14.660 & 0.035 \\
z1\_ng4 & 0.435 & 156.1 &    70 & 14.910 & 0.061 \\
z1\_ng5 & 0.417 & 237.4 &    18 & 15.191 & 0.070 \\
z2\_ng0 & 0.627 &  30.1 & 3{,}488 & 14.088 & 0.039 \\
z2\_ng1 & 0.624 &  45.7 & 1{,}384 & 14.203 & 0.056 \\
z2\_ng2 & 0.627 &  68.2 &   436 & 14.332 & 0.076 \\
z2\_ng3 & 0.620 & 102.1 &   138 & 14.615 & 0.114 \\
z2\_ng4 & 0.608 & 154.1 &    30 & 14.800 & 0.160 \\
\rowcolor{red!10} z2\_ng5 & 0.646 & 227.0 & 2 & 15.942 & 0.014 \\
\bottomrule
\end{tabular}
\caption{NFW + miscentering posterior means per stack. The
\textcolor{red!70!black}{red} row is the $N=2$ outlier whose fit hits prior
edges ($C\!\to\!15$, $F_{\rm miss}\!\to\!0$); it is excluded from the
power-law fits below. Under our new convention (see~\autoref{sec:caveats})
\texttt{z0\_ng5} ($N=6$) would also be excluded going forward.}
\label{tab:fits18}
\end{table}

\begin{figure}[h]
\centering
\includegraphics[width=\linewidth]{mr_profiles.png}
\caption{Stacked $\DS$ profiles for the 18 $(z, N_{\rm gals})$ cells.
One panel per redshift slice, six richness curves overlaid per panel
(viridis gradient, low $\lambda$ → blue). $\DS$ amplitude scales smoothly
with both richness (at fixed $z$) and inversely with redshift (at fixed
$\lambda$). Log--log axes.}
\label{fig:mrprof}
\end{figure}

\begin{figure}[h]
\centering
\includegraphics[width=\linewidth]{mr_overlay.png}
\caption{Best-fit miscentered NFW models (solid curves) overlaid on the
stacked profiles. The model captures both the $R^{-1}$ inner behaviour
and the high-$R$ turnover where the 2-halo term becomes relevant.}
\label{fig:mroverlay}
\end{figure}

\begin{figure}[h]
\centering
\includegraphics[width=0.7\linewidth]{mr_relation.png}
\caption{WaZP mass--richness relation. Filled circles are the 17 stacks
used in the fit; the open marker is the \texttt{z2\_ng5} ($N=2$) outlier
excluded for its prior-edge collisions. Coloured dashed lines are per-z
1D fits at $\lambda_{\rm piv}=40$; the solid black line is the joint
$M(\lambda, 1+z)$ relation evaluated at $z=z_{\rm piv}=0.35$.}
\label{fig:mrrel}
\end{figure}

\paragraph{Power-law M--$\lambda$ fits.}
Excluding \verb|z2_ng5|, we fit two models by weighted least squares on
$\log_{10}\Mtwoh$ vs $\log_{10}\lambda$, with weights
$1/\sigma_{\log M}^2$ from the posterior $16/84$ half-width.

\emph{Per-$z$ 1D fits} at $\lambda_{\rm piv}=40$:
\begin{align*}
\langle z\rangle = 0.274:\quad
\log_{10}\Mtwoh &= 14.374\pm0.015 + (1.074\pm0.058)\,\log_{10}(\lambda/40) \\
\langle z\rangle = 0.440:\quad
\log_{10}\Mtwoh &= 14.286\pm0.014 + (1.012\pm0.066)\,\log_{10}(\lambda/40) \\
\langle z\rangle = 0.621:\quad
\log_{10}\Mtwoh &= 14.179\pm0.013 + (0.761\pm0.095)\,\log_{10}(\lambda/40)
\end{align*}
The lowest two redshift bins are mutually consistent in slope
($\alpha\!\approx\!1.0$); the highest-$z$ bin sits about $0.2\Msun$ lower at
fixed $\lambda$ and prefers a shallower slope, but covers a narrower
$\lambda$ range and has the smallest stack populations.

\emph{Joint fit} with explicit z evolution
($\lambda_{\rm piv}=40$, $z_{\rm piv}=0.35$):
\[
  \log_{10}\Mtwoh = (14.335\pm 0.011)
                  + (1.030\pm 0.035)\,\log_{10}(\lambda/40)
                  + (-1.728\pm 0.256)\,\log_{10}((1+z)/1.35).
\]
The slope $\alpha\approx 1.03$ agrees well with the literature
(McClintock+18 reports $\sim 1.36$ for DES Y1 redMaPPer at $z\!\sim\!0.4$;
WaZP being a photometric finder with a different richness definition is
expected to differ). The redshift slope $\beta\approx -1.7$ is a strong,
\emph{negative} z-evolution -- i.e.\ at fixed $\lambda$, masses drop with
redshift. This is partly degenerate with the underestimation of high-$z$
sources (DNF photo-zs broaden at higher $z$, biasing $\Sigma_{\rm crit}$ and
hence $\DS$), so the magnitude of $\beta$ should not be over-interpreted at
this stage.

\paragraph{Outlier handling.} An earlier version of the joint fit
\emph{included} \verb|z2_ng5|, which gave $\alpha=+1.68\pm0.03$ and
$\beta=+3.09\pm0.19$. The factor-of-three swing on $\beta$ between
``include'' and ``exclude'' is itself a robust diagnostic that the highest-z,
highest-$\lambda$ stack is dominated by a single object's likelihood pinning
the prior edge. We remove it from the published M--$\lambda$ relation and
mark it on the plot with an open marker.

\section{$\dm$-dependence of the M--$\lambda$ relation
(\texttt{wazp\_mass\_richness\_by\_delta\_m})}
\label{sec:dmsplit}

To test whether the magnitude gap is a second observable of the
mass--richness relation, we repeat \autoref{sec:massrichness} after splitting
each $(z, N_{\rm gals})$ cell in two halves at the cell median $\dm$
(``low $\dm$'' / ``high $\dm$''). Clusters with non-finite or non-positive
$\dm$ are dropped first; $\dm\le 0$ values are likely cluster-finder
artefacts (BCG mis-assignment, the brightest member residing closer to the
centre than the catalog BCG). The $\dm>0$ cut drops $142$ clusters, leaving
$15{,}825$.

\paragraph{Bin populations.} Per-cell median $\dm$ ranges from
$\sim 0.46$ to $\sim 0.99\,\mathrm{mag}$ across the 9 cells (lower in
high-$z$, mostly insensitive to richness). Cell populations balance
symmetrically by construction (50/50 within each cell). The largest cells
($z\in[0.35,0.55)$, $N_{\rm gals}\in[25,38)$) contain
$\sim 2{,}268$ clusters per $\dm$ half; the smallest cells
($z>0.55$, $N_{\rm gals}>200$) only $1$ per half and are dropped from the
fit.

\paragraph{Fit results.} The 36-stack pool completed with 34 successful
fits and 2 failures ($\mathtt{z0\_ng5\_dm0}$ at $N=3$ and
$\mathtt{z2\_ng5\_dm0}$ at $N=1$ -- both below the cell-population threshold
where nautilus' bound splitter can compute a 4D ellipsoid). The retry logic
(\texttt{n\_live} bumped to $\max(4n_{\rm live}, 2000)$,
$\texttt{n\_points\_min}=50$, $\texttt{split\_threshold}=400$) caught most of
the marginal stacks. The fits-out pickle survives both failures by storing a
\texttt{failed: True} marker, and the downstream M--$\lambda$ regression
filters those out.

\paragraph{Per-$\dm$-group M--$\lambda$ relation.}
Excluding failed fits, weighted least squares on the surviving stacks gives
the two power-law relations
\begin{align*}
\text{low}\ \dm\ \ (N_{\rm stacks}=16):\quad
\log_{10} M &= (14.328\pm 0.017) + (1.063\pm 0.039)\log_{10}(\lambda/40) \\
                  &\quad + (-1.929\pm 0.383)\log_{10}((1+z)/1.35), \\[2pt]
\text{high}\ \dm\ (N_{\rm stacks}=17):\quad
\log_{10} M &= (14.359\pm 0.013) + (0.935\pm 0.023)\log_{10}(\lambda/40) \\
                  &\quad + (-1.370\pm 0.323)\log_{10}((1+z)/1.35).
\end{align*}
The high-minus-low difference per parameter is
\begin{center}\small
\begin{tabular}{lr}
\toprule
parameter & high $-$ low \\
\midrule
$\Delta\alpha$ (slope $\lambda$)  & $-0.128 \pm 0.045$ \quad \textbf{($-2.85\sigma$)} \\
$\Delta\beta$ (slope $1+z$)       & $+0.559 \pm 0.501$ \quad $(+1.12\sigma)$ \\
$\Delta\log M_0$                  & $+0.031 \pm 0.021$ \quad $(+1.45\sigma)$ \\
\bottomrule
\end{tabular}
\end{center}

The $\lambda$ slope is significantly shallower in the high-$\dm$ sample
($-2.85\sigma$). The normalisation $M_0$ at $\lambda_{\rm piv}=40$,
$z_{\rm piv}=0.35$ is consistent within $1.5\sigma$. Translated to a
physical statement: \emph{at fixed low-to-moderate richness
($\lambda\lesssim 40$), high-$\dm$ clusters lie at the same or slightly
higher mass than low-$\dm$ ones; at high richness ($\lambda\gtrsim 100$)
the high-$\dm$ relation predicts lower mass than the low-$\dm$ relation.}
This is the directional sense expected if high-$\dm$ systems are more
relaxed/older and have therefore had more time for their satellite
luminosity functions to develop a gap -- but this interpretation is
preliminary at present statistical power and given the photometric-$z$
caveats listed in \autoref{sec:caveats}.

\begin{figure}[h]
\centering
\includegraphics[width=0.75\linewidth]{mr_dm_contours.png}
\caption{Posterior contours of the M--$\lambda$ power-law parameters
$(\alpha, \beta, \log M_0)$ for the low-$\dm$ sample (blue) and high-$\dm$
sample (red). The 1- and 2-$\sigma$ levels are shown. The two posteriors
overlap in $\beta$ and $\log M_0$; they separate cleanly in $\alpha$, with
the high-$\dm$ slope being $\Delta\alpha = -0.128\pm 0.045$ shallower
($-2.85\sigma$). Drawn from $20{,}000$ samples per group of the analytic
Gaussian posterior.}
\label{fig:contours}
\end{figure}

\section{Technical issues encountered and fixes}

\paragraph{\texttt{destacking.data.magnitude\_gaps}.}
The function was unrunnable in the repository: missing colon on the
\texttt{def} line, missing \texttt{numpy} / \texttt{astropy.table} imports,
and four typos using \texttt{centre\_galaxies['key']} where
\texttt{centre\_members['key']} was meant. The variance term used
\texttt{delta**2} instead of \texttt{mag\_gap**2}. We fixed these minimally
in place (no algorithmic change), added an early
\texttt{(nan, nan)} return when no satellites fall inside $r_{\rm cut}$,
and smoke-tested the result.

\paragraph{Photo-z source.} The DES Y3 metacal HDF5 has no source-redshift
column. The DNF file \verb|DESY3_GOLD_2_2.1_DNF.h5| is row-aligned with
metacal (same $N$, same ordering of \verb|coadd_object_id|). We use
\verb|zmean_sof|, validate row alignment at the first and last index, and
keep only the row-aligned column to save memory.

\paragraph{Nautilus bound failure.} For some stacks the nautilus bound
splitter raised
\texttt{ValueError("Number of points must be larger than number dimensions")}.
Cause: the multi-ellipsoid Union splitter produced a sub-cluster with
fewer than $D+1=5$ points, which is undefined for an ellipsoid in 4D. Fix:
\texttt{fit\_one} now retries with $n_{\rm live} \to \max(4n_{\rm live},
2000)$, $n_{\rm points\_min}=50$, $\texttt{split\_threshold}=400$. A second
failure is reported as a structured marker rather than crashing the pool,
and the loop writes \texttt{FITS\_PKL} after every completed fit so
interruptions cost at most one stack.

\paragraph{$\DS$ plots.}
The convention is log--log. Linear-$y$ crushes the inner profile (where
the signal carries most of the concentration/miscentering information) into
the bottom of the panel. All $\DS$ vs $R$ plots in
\autoref{sec:percluster}--\ref{sec:dmsplit} use
\texttt{set\_xscale('log')} \emph{and} \texttt{set\_yscale('log')}.

\section{What the numbers say}

\begin{itemize}
\item The WaZP real-data sample of $15{,}967$ clusters carries a full set of
position, redshift, richness, magnitude-gap, and (post-pipeline) $\DS$ vector,
making it a self-contained analysis catalog.
\item The mean of magnitude gaps ($\sim 0.62$\,mag in $i$) is consistent with
a population dominated by hierarchical assembly, with a long tail at high
$\dm$ that may pick out relaxed/fossil systems.
\item At fixed $\lambda$, masses at $z\sim 0.27$ and $z\sim 0.44$ are
mutually consistent and follow a near-unit slope
($\alpha=1.07\pm0.06$ and $1.01\pm0.07$ respectively).
\item The highest-redshift bin prefers a shallower $\lambda$ slope
($\alpha=0.76\pm0.10$) and lower normalization, but is the most prone to
photo-$z$-driven biases on $\Sigma_{\rm crit}$.
\item The joint fit gives $\alpha = 1.03\pm0.04$ and a redshift slope
$\beta = -1.7\pm0.3$. The latter changes sign and magnitude depending on
whether the $N=2$ outlier (\texttt{z2\_ng5}) is included, so we treat it as
preliminary.
\item Whether $\dm$ is a second observable of cluster mass at fixed
$\lambda$ is the explicit question being answered by the companion
notebook; the analysis is ready and awaits the 36-stack fit run.
\end{itemize}

\section{Caveats and follow-ups}
\label{sec:caveats}

\begin{enumerate}
\item Source photo-$z$s are point estimates (\verb|zmean_sof|); for a publication
result we would integrate over the per-source PDF (or at minimum the DNF
$z_{\rm MC}$ Monte Carlo draw) and propagate the $\Sigma_{\rm crit}$
uncertainty into the stack covariance.
\item No boost-factor or $1+\Sigma_{\rm crit}\,e_t$-shear correction beyond
what is internal to \texttt{xlensing.data.metacal\_cluster\_lensing} is
applied.
\item The four NFW parameters $\{\Mtwoh, C, F_{\rm miss}, \sigma_{\rm off}\}$
are heavily degenerate at the radial range used; tighter priors on
miscentering from simulations would sharpen $\Mtwoh$.
\item The cosmology is \verb|Planck18| throughout; uncertainty on
$D_A(z)$ has not been propagated.
\item For the highest-$\lambda$ stacks ($N\lesssim 30$) the bootstrap
covariance is noisy and the fits become prior-dominated; future work should
either widen $\lambda$ bins above $\sim 100$ or marginalize over a noise
floor.
\item \emph{Convention going forward}: stacks with fewer than $10$ clusters
should be computed for completeness but excluded from \emph{all} downstream
relations (per-z 1D fits, joint $M(\lambda, z)$ fit, $\dm$-split
comparison). Under this rule, the present 18-stack relation drops
\texttt{z0\_ng5} ($N=6$) and \texttt{z2\_ng5} ($N=2$) -- only
\texttt{z2\_ng5} is dropped in the report's numbers. The
$\dm$-split relation drops both \texttt{z0\_ng5\_dm[01]} ($N=3$ each),
\texttt{z2\_ng5\_dm[01]} ($N=1$ each, already failed), and
\texttt{z1\_ng5\_dm[01]} ($N=9$ each).
\end{enumerate}

\end{document}
